% Modeling strategy section.
% Requires: \usepackage{amsmath, amssymb}

\section{Modeling strategy}

\subsection{The outcome}

The dependent variable is the number of automated license plate readers
recorded in a Census tract. Across the 84{,}415 tracts in the contiguous United
States the count has mean $1.47$ and variance $7.92$, and $54.0$ percent of
tracts contain none. Conditional on containing at least one, the mean is $3.20$
and the variance $11.67$.

Two features of this distribution constrain the choice of model. The first is
overdispersion: the variance-to-mean ratio is $5.37$ unconditionally, and
remains $3.64$ among tracts with at least one camera, so the excess dispersion
is not merely an artefact of the mass at zero. The second is the size of that
mass. A Poisson distribution with the observed mean implies a zero share of
$22.9$ percent against the $54.0$ percent observed, a discrepancy of thirty-one
percentage points that no adjustment to the mean function can absorb.

\subsection{Why a hurdle rather than a zero-inflated model}

Both hurdle and zero-inflated specifications accommodate more zeros than a
single count process would generate, but they embody different substantive
claims. A zero-inflated model posits two kinds of zero --- structural zeros
that could never have been positive, mixed with sampling zeros drawn from the
count process. A hurdle posits a single process governing whether the count
crosses zero, and a second, separate process governing its magnitude once it
has \citep{Mullahy1986}.

The deployment of license plate readers matches the hurdle. An agency first
decides whether to place any reader in a given area, and then how many; there
is no natural class of tract that could never receive a camera, since almost
any tract with a road could. We refer to these two decisions as the
\emph{adoption} and \emph{intensity} margins.

The distinction is not merely conceptual. In our data the poverty rate enters
the two margins with opposite signs --- negative on adoption, positive on
intensity. A single count model averages these into approximate zero and would
support the conclusion that poverty is unrelated to surveillance. The hurdle
recovers both.

We note that the two specifications are, on these data, statistically
indistinguishable: a zero-inflated negative binomial and a hurdle negative
binomial fit with the same covariates differ by two AIC points on $82{,}916$
observations. The choice between them therefore rests on the substantive
argument above rather than on fit.

\subsection{Specification}

Let $Y_i$ denote the camera count in tract $i$, located in county $c(i)$. Write
$\pi_i$ for the probability that $Y_i$ is positive and $\mu_i$ for the
conditional mean of the underlying count process. The hurdle density is

\begin{equation}
\Pr(Y_i = y) =
\begin{cases}
1 - \pi_i, & y = 0,\\[6pt]
\pi_i \, \dfrac{f_{\mathrm{NB}}(y \mid \mu_i, \theta)}
               {1 - f_{\mathrm{NB}}(0 \mid \mu_i, \theta)}, & y \geq 1,
\end{cases}
\label{eq:hurdle}
\end{equation}

where $f_{\mathrm{NB}}$ is the negative binomial (NB2) mass function. The
adoption margin is a logit,
\begin{equation}
\operatorname{logit}(\pi_i) = \alpha + \mathbf{x}_i'\boldsymbol{\beta} + u_{c(i)},
\end{equation}
and the intensity margin a zero-truncated NB2,
\begin{equation}
\log(\mu_i) = \gamma + \mathbf{x}_i'\boldsymbol{\delta} + v_{c(i)}.
\end{equation}

The two margins share the covariate vector but not its coefficients:
$\boldsymbol{\beta} \neq \boldsymbol{\delta}$ is the object of interest, since a
covariate may govern whether an area is surveilled at all without governing how
heavily. Because the hurdle log-likelihood separates additively into a term in
$\boldsymbol{\beta}$ and a term in $(\boldsymbol{\delta}, \theta)$, the two
margins may be estimated separately without loss of efficiency.

\subsection{Estimators considered and rejected}

Three natural approaches fail, each for a different reason, and the failures
motivate the estimator we adopt.

\paragraph{Poisson.} Rejected on the evidence above. Fitted with our full
covariate set the Poisson hurdle returns an AIC of $276{,}879$ against
$251{,}212$ for the negative binomial equivalent. Its standard errors would be
understated by roughly a factor of $\sqrt{5.37} \approx 2.3$.

\paragraph{Standard hurdle software without geographic controls.} The
conventional implementation \citep[\texttt{pscl};][]{Zeileis2008} estimates
both margins by maximising a likelihood over a dense design matrix. Absorbing
county effects would require some $2{,}270$ indicator variables over $82{,}728$
observations, and unlike the linear case there is no within-transformation that
sweeps them out of a nonlinear model. Estimated without county controls, the
coefficient on share Black is $0.989$; with them it is $0.495$. Half of the
apparent association is between-county sorting, so leaving geography
uncontrolled is not a viable option.

\paragraph{Fixed-effects count software applied to the positive subsample.} A
tempting shortcut is to estimate the intensity margin by applying a
fixed-effects negative binomial to the tracts with $Y_i > 0$. This is a
misspecification rather than an approximation. The second margin of a hurdle is
a \emph{zero-truncated} distribution, and fitting an untruncated negative
binomial to truncated data attenuates every slope by roughly half and biases
the dispersion parameter by a factor of five in our data (for example, the
coefficient on renter-occupied share is $0.93$ untruncated against $1.42$
correctly truncated; $\theta$ is $2.79$ against $0.53$). No available
fixed-effects estimator implements a zero-truncated negative binomial family.

\subsection{Estimation: correlated random effects}

We therefore estimate both margins by maximum likelihood with a county random
intercept, augmented with the county mean of every covariate --- the
correlated-random-effects device of \citet{Mundlak1978}. Writing the county
effect as a function of the covariate means,
$\alpha_c = \bar{\mathbf{x}}_c'\boldsymbol{\gamma} + \omega_c$, and substituting
into the model, the coefficient on $\mathbf{x}_i$ becomes the within-county
effect. The random-effects assumption is thereby relaxed to hold conditional on
the county means, and the estimand is that of fixed effects while the
likelihood remains one that admits zero truncation.

We verify the device rather than assume it. On the adoption margin, where
conventional fixed effects are estimable, the two coefficient vectors correlate
at $0.9999$ with a median absolute difference of $3.0$ percent. The
correlated-random-effects estimator additionally retains the $1{,}092$ counties
($4{,}656$ tracts) that fixed effects must discard for exhibiting no
within-county variation in the outcome.

In reporting, the plain covariates are interpreted as within-county effects;
the county-mean terms are nuisance parameters absorbing between-county
variation and are estimated but not tabulated.

\subsection{Interpretation}

We treat the estimates as descriptive rather than causal. Cameras are sited by
law enforcement agencies and private operators, not by tracts; tract
characteristics proxy for what a decision-maker observed when siting, and the
agency rather than the tract is the decision unit. The coefficients answer a
conditional descriptive question: within a county, what kinds of tracts contain
cameras, and how many?

Two limitations bear directly on that reading. First, the outcome records
cameras that volunteers have \emph{mapped}, not cameras that exist; because
mapping effort plausibly rises with income and education, coefficients on those
covariates are partly a discovery artefact. Second, roughly fourteen percent of
the cameras whose operator is identifiable belong to retail chains, whose
siting is a corporate decision unrelated to local policing; where operator
information permits, we report estimates separately by operator type.
